%% outlook.tex
\section{Outlook}\label{sec:outlook}

The worked examples used standard free energies, turnover numbers, and
regulatory models to set the factors in Equation~\eqref{eq:general-bound}.
Time-course data can instead be used to estimate turnover numbers, saturation
constants, allosteric factors $\theta_j$, and expression controls $u_j$ and
$w_j$. Tabulated values remain useful as prior information when dynamic data
are unavailable. The same approach can be applied reaction by reaction in
genome-scale models.
These models have many flux distributions that satisfy
$\mathbf S\hat{\mathbf v}=\mathbf 0$, so condition-specific bounds are
especially important. Setting an unknown factor to one preserves feasibility
but may leave many fluxes unresolved. Time-course data also allow models to
retain accumulation and growth dilution.
Cell-free systems have $\mu=0$ but can have substantial accumulation. Growing
cell cultures can have both terms. When measurements resolve these changes,
fluxes should be estimated from Equation~\eqref{eq:dilution-constraint}
instead of imposing Equation~\eqref{eq:sv0}.
